\section{Download Speed}
\label{sec:cses-download-speed}

\problemstatement{Given a directed graph of $n$ nodes and $m$ edges, each
with a capacity (bandwidth), find the maximum data rate from node $1$
(source) to node $n$ (sink).}

\begin{patternbox}
Maximum throughput under capacities is the \textbf{maximum flow} problem.
\textbf{Dinic's algorithm} computes it efficiently: repeatedly build a
level graph with BFS, then push blocking flow along shortest augmenting
paths with DFS.
\end{patternbox}

\subsection*{Level graphs and blocking flow}

A flow assigns each edge a value up to its capacity, with flow conserved at
every node except source and sink; the max flow is the greatest total
leaving the source. Dinic alternates two phases: a BFS labels each node with
its distance (level) from the source in the residual graph; a DFS then sends
as much flow as possible along edges that go strictly one level deeper (the
blocking flow). Reverse edges are maintained so earlier flow can be
rerouted. When BFS can no longer reach the sink, the flow is maximum.

\begin{ideabox}[Augment Along Shortest Paths, in Bulk]
Restricting augmenting paths to the level graph (only forward-level edges)
and saturating them all before rebuilding gives Dinic its speed: the number
of BFS phases is $O(V)$, far fewer than pushing one path at a time.
\end{ideabox}

\subsection*{Algorithm}

\begin{enumerate}
  \item Build the flow network with forward capacities and zero-capacity
        reverse edges.
  \item While BFS reaches the sink: DFS blocking flow along level edges,
        accumulating flow.
  \item Output the total flow.
\end{enumerate}

\subsection*{Dry run}

\dryrun{$1\!\to\!2(3)$, $1\!\to\!3(2)$, $2\!\to\!4(2)$, $3\!\to\!4(3)$;
sink $4$.}

\begin{itemize}
  \item Push $2$ along $1\!\to\!2\!\to\!4$ and $2$ along
        $1\!\to\!3\!\to\!4$.
  \item Max flow $=4$ (limited by outgoing/incoming capacities).
\end{itemize}

\subsection*{C++ implementation}

\begin{lstlisting}
#include <bits/stdc++.h>
using namespace std;

struct Dinic {
    struct Edge { int to; long long cap; int rev; };
    vector<vector<Edge>> g;
    vector<int> level, it;
    Dinic(int n) : g(n), level(n), it(n) {}
    void addEdge(int u, int v, long long c) {
        g[u].push_back({v, c, (int)g[v].size()});
        g[v].push_back({u, 0, (int)g[u].size() - 1});   // reverse edge
    }
    bool bfs(int s, int t) {
        fill(level.begin(), level.end(), -1);
        queue<int> q; level[s] = 0; q.push(s);
        while (!q.empty()) {
            int u = q.front(); q.pop();
            for (auto& e : g[u])
                if (e.cap > 0 && level[e.to] < 0) { level[e.to] = level[u] + 1; q.push(e.to); }
        }
        return level[t] >= 0;
    }
    long long dfs(int u, int t, long long f) {
        if (u == t) return f;
        for (int& i = it[u]; i < (int)g[u].size(); i++) {
            Edge& e = g[u][i];
            if (e.cap > 0 && level[u] < level[e.to]) {
                long long d = dfs(e.to, t, min(f, e.cap));
                if (d > 0) { e.cap -= d; g[e.to][e.rev].cap += d; return d; }
            }
        }
        return 0;
    }
    long long maxflow(int s, int t) {
        long long flow = 0;
        while (bfs(s, t)) {
            fill(it.begin(), it.end(), 0);
            long long f;
            while ((f = dfs(s, t, LLONG_MAX)) > 0) flow += f;
        }
        return flow;
    }
};

int main() {
    int n, m;
    cin >> n >> m;
    Dinic dinic(n + 1);
    for (int i = 0; i < m; i++) {
        int a, b; long long c; cin >> a >> b >> c;
        dinic.addEdge(a, b, c);
    }
    cout << dinic.maxflow(1, n) << "\n";
    return 0;
}
\end{lstlisting}

\begin{complexitybox}
\textbf{Time Complexity.} $O(V^2 E)$ worst case (Dinic), far faster in
practice. \textbf{Space Complexity.} $O(V+E)$.
\end{complexitybox}

\begin{takeawaybox}
Maximum throughput is maximum flow; Dinic's level-graph blocking flow is the
standard efficient solver. This one algorithm is the engine behind the next
three problems -- min cut, matching, and edge-disjoint paths -- each just a
different network.
\end{takeawaybox}
